% Independent geometric input for D.2.44.  No separate theorem ID is used.
% Saved for root/fourfolds audit; the owning D.2.44 statement supplies the
% actual homology and critical-cycle computations.
\subsubsection*{From a critical-cycle linking number to the relative three-frame}

The following argument applies to the semistable primitive-Mumford rows
under consideration.  Let $X$ be their simply connected complex
threefold, with
\[
 H_2(X;\mathbb Z)=H_3(X;\mathbb Z)=\mathbb Z/c,
 \qquad c_1(X)=c_2(X)=0,
 \qquad c\text{ even}.
\]
The asserted equality for $H_3$ is an equality of homology groups.
Let $T=\langle e_1,e_2,\gamma\rangle$ be the regular three-torus,
with the normal frame supplied by the complementary circle and the
two oriented base directions.  Assume, as independently established
by the marked cap calculation, that $[T]$ generates $H_3(X)$.
The conclusion of this input is
\[
 \overline\gamma(X)=\frac c2
       \bigl(\operatorname{lk}(T,C_{\rm crit})\bmod2\bigr)
       \quad\text{in }H^4(X;\mathbb Z)=\mathbb Z/c.
 \tag{D.F1}
\]
Here $C_{\rm crit}$ is the literal complex critical-curve cycle, and
the integer linking number uses its integral nullhomology.  The sign
of the chosen generator does not affect (D.F1).  The proof below
uses Zhubr's proved Theorem~1, not his general Conjecture~1.

\paragraph{The actual vertical complex and its local obstruction.}
Write $L$ for the target line of the differential after removing any
common multiple-fibre factor.  In the rows in question the toric
semistable models are reduced.  The kernel $V$ is an abstract locally
free complex rank-two sheaf, although its map into $TX$ drops rank on
the critical locus.  For $q=xy$, local generators are
$(x,-y,0)$ and $(0,0,1)$; for $q=xyz$, they are
$(x,-y,0)$ and $(x,0,-z)$.  Thus the exact sheaf complex is
\[
 0\longrightarrow V\longrightarrow TX\xrightarrow{df}L
 \longrightarrow L\otimes\mathcal O_{\rm Crit}\longrightarrow0.
 \tag{D.F2}
\]
The scalar circle direction $\delta$ is a nowhere-zero section of
the abstract bundle $V$.  In these logarithmic generators it is a
constant nonzero vector; its image in $TX$ may vanish on a fixed
curve.  Since $H^2(X;\mathbb Z)=0$, both $V/\mathbb C\delta$ and
$L$ are topologically trivial.  Consequently
$V\oplus L$ has a global complex frame whose first vertical vector
is $\delta$.

The virtual difference in (D.F2) is
$[TX]-[V]-[L]=-[L\otimes\mathcal O_{\rm Crit}]$.
At a generic point of each critical curve it has the node model.
On the normal three-sphere of $xy=0$, the two columns
\[
 \frac1{\sqrt{|x|^2+|y|^2}}
   \begin{pmatrix}x&\overline y\\-y&\overline x\end{pmatrix}
 \tag{D.F3}
\]
give the unit $SU(2)$ clutch.  Its second-Chern obstruction is a
generator, up to the common orientation convention.  Every complex
critical curve therefore occurs with coefficient one.  The
codimension-three triple points do not add a codimension-two
component.  In particular $c_2(TX)=\operatorname{PD}(C_{\rm crit})$.
This is a relative obstruction statement as well as an equality of
absolute Chern classes: outside $C_{\rm crit}$, a splitting of
(D.F2) identifies $TX$ with $V\oplus L$, and its degree-four
extension obstruction is exactly that cycle.

Choose a point away from the torus, the critical cycle and its
bounding chain.  The bundle $TX$ on $X\setminus\{\mathrm{pt}\}$
has a complex frame $f$: on a five-dimensional spine, a rank-three
complex bundle with $c_1=c_2=0$ is trivial, since the only relevant
nonzero homotopy group of $BSU(3)$ is
$\pi_4BSU(3)=\mathbb Z$.
The determinant of the exact sequence extends across its
codimension-two exceptional set.  Choose the determinants of the
two global frames compatibly with this isomorphism.  Comparing
$f$ with the frame of $V\oplus L$ through the splitting on the
regular locus now gives a map to $SU(3)$.
If $k$ is its degree-three class evaluated on $T$, (D.F3) gives
\[
                   k=\pm\operatorname{lk}(T,C_{\rm crit}).
 \tag{D.F4}
\]
Indeed its connecting image in relative degree four is the
critical-cycle class.  Changing a global frame changes its
evaluation by the pairing of an integral class of $H^3(X)$ with
the torsion cycle $[T]$, which is zero.  This also proves that the
integer in (D.F4) is independent of the chosen bounding chain.

The $H^3$ part of this comparison is the same as that with the
actual constant complex frame on the regular fibre.  For the
vertical factor, both frames have first vector $\delta$, so their
transition lies in the stabilizer of this vector, an upper
triangular group retracting to $\mathbb C^*$.  The target-line
transition also lies in $\mathbb C^*$, and changing the horizontal
lift adds only a contractible upper triangular factor.  None of
these changes has a degree-three class.  This is why the complete
critical cycle, rather than an arbitrary half of its walls,
calculates the flat-versus-global framing degree.

\paragraph{An embedded normal-framed conversion to a sphere.}
Here is the spin-sensitive geometric step.  More generally, let
$N^3$ be a connected oriented normally framed submanifold of a
simply connected framed six-manifold $U$.  There is a normally
framed embedded cobordism in $U\times[0,1]$ from $N$ to an
embedded $S^3\subset U\times\{1\}$, and it preserves the stable
tangent-framing bordism class induced from the ambient frame.

To prove this, use the spin structure on $N$ induced by the ambient
frame and the specified normal frame.  Kaplan's constructive
Theorem~3.1, with Definition~1.7 and Remark~1.8, produces a compact
four-manifold with one zero-handle and only two-handles, whose
boundary has this specified almost framing
\cite{mapsKaplan79}.  Its spin structure therefore restricts to
the specified one.  Delete the zero-handle and reverse the
cobordism.  The result is a spin cobordism from $N$ to $S^3$
consisting relative to $N$ only of two-handles.

Realize these handles as ambient surgeries in $U$ itself.
Each attaching circle bounds a disk because $\pi_1(U)=0$.
Prescribe its inward boundary collar to point along the first
normal vector of $N$.  General position makes the disk embedded,
its interior disjoint from $N$, and all the disks pairwise
disjoint: the relevant strict inequalities are $2+3<6$ and
$2+2<6$.  Along its boundary, the rank-four normal bundle of the
disk is framed by the two specified tangent directions of the
two-handle attachment and the remaining two normal vectors of
$N$.  Its extension obstruction is in $\pi_1SO(4)=\mathbb Z/2$.
It is precisely the obstruction to extending the induced spin
structure over that two-handle, and hence vanishes.  Extend this
four-frame over the disk and thicken the first two directions to
the prescribed embedded $D^2\times D^2$ surgery block.  The
remaining directions, together with the outward direction at
its new boundary, give the new normal three-frame.

These ambient surgeries give the required $S^3$ in $U$, not just
a sphere somewhere in a seven-dimensional product.  Their
traces give the asserted embedded cobordism.  Equivalently, the
relative normal-framing obstruction of the trace is its relative
$w_2$, which vanishes by the chosen spin cobordism; there are no
higher relative obstructions because there are only two-handles.
The ambient frame and this normal frame induce a stable tangent
frame throughout the trace.  Thus
\[
       [N]_f=[S^3]_f\quad\text{in }\Omega^{\rm fr}_3=\mathbb Z/24,
       \qquad [S^3]=[N]\quad\text{in }H_3(U;\mathbb Z).
 \tag{D.F5}
\]
Apply this to $U=X\setminus\{\mathrm{pt}\}$ and $N=T$.
The resulting sphere is an actual torsion generator.
If $c=2^r m$ with $m$ odd, take the connected sum of $m$
disjoint parallel copies, using disjoint joining arcs in $U$.
This is again an embedded normally framed three-sphere and
represents the generator $m[T]$ of the two-primary subgroup.
Its framed class is multiplied by the odd number $m$, so its
parity is unchanged.  In the seven semistable rows $m=1$.
Kaplan's Example~3.2 explicitly distinguishes the eight almost
framings of $T^3$: seven have an index-zero filling and the
remaining one has an index-eight filling.  Therefore this use of
his theorem does not silently replace the Lie spin structure by
the spin structure of a zero-framed Borromean trace.

\paragraph{The parity retained by this conversion.}
The flat tangent and normal frames induce the Lie frame on
$T^3$, whose stable framed class is $\eta^3=12\nu$.
For a transition map $g:T^3\to U(3)$, split off
$\operatorname{diag}(\det g,1,1)$.  The remaining $SU(3)$ map is,
on a three-dimensional complex, homotopic to a degree-$k$ map
supported in a ball.  Realification sends its $\pi_3$ generator
to a stable $SO$ generator: a unit $c_2$ has underlying real
$p_1=\mp2$.  Its effect on the framed class is therefore
$\pm k\nu$.

The determinant factor cannot change its parity.  A map
$T^3\to S^1$ is a character up to homotopy; after an integral
change of torus basis it depends on just one circle.  Stably
conjugate its rotation plane to that circle plus one trivial
normal direction.  It then changes just the framing of that
circle.  The resulting product-framed torus represents either
$0$ or $\eta^3=12\nu$, both even multiples of $\nu$.
Consequently (D.F4)--(D.F5) give
\[
                   [S^3]_f\bmod2=k\bmod2.
 \tag{D.F6}
\]
Changes of the sphere's normal three-frame add the image of
$\pi_3SO(3)$, which is twice the stable generator, and do not
change (D.F6).

\paragraph{The universal rank-three complement in the trivial ambient bundle.}
For completeness the exact obstruction comparison needed here
has a direct universal proof.  Set
$V=V_{6,3}=\operatorname{Spin}(6)/\operatorname{Spin}(3)$,
and let $\iota\in H^3(V;\mathbb Z/2)$ be its nonzero class.
Zhubr's Lemma~1 gives
\[
 \pi_1V=\pi_2V=\pi_4V=0,\qquad \pi_3V=\mathbb Z/2,
 \qquad
 \pi_3\operatorname{Spin}(3)\longrightarrow
 \pi_3\operatorname{Spin}(6): n\longmapsto2n.
\]
Hurewicz gives $H_3(V;\mathbb Z)=\mathbb Z/2$.
For the next degree, map the Moore space
$M(\mathbb Z/2,3)=S^3\cup_2D^4$ to $V$ by a generator of
$\pi_3V$ and a nullhomotopy of its double.  This map induces
an isomorphism on $\pi_3$ and a surjection on $\pi_4=0$,
so its mapping-cylinder pair is four-connected.  Relative
Hurewicz, together with $H_4(M(\mathbb Z/2,3);\mathbb Z)=0$,
therefore gives $H_4(V;\mathbb Z)=0$.
Thus $H^4(V;\mathbb Z)=\mathbb Z/2$, generated by $\beta\iota$.
The universal orthogonal three-plane is spin.  In the Serre
sequence of the displayed principal $\operatorname{Spin}(3)$
bundle, the low-degree edge and transgression sequence is
\[
 H^3(\operatorname{Spin}(6);\mathbb Z)=\mathbb Z
 \xrightarrow{\,2\,}
 H^3(\operatorname{Spin}(3);\mathbb Z)=\mathbb Z
 \xrightarrow{d_4}H^4(V;\mathbb Z)=\mathbb Z/2.
 \tag{D.F7}
\]
The transgression of the fibre generator is exactly the first
obstruction to framing the spin three-plane.  It is nonzero by
(D.F7), and therefore equals $\beta\iota$.
For any map $v:U\to V$ in a globally trivial oriented rank-six
bundle this proves
\[
                         o_4(v^\perp)=\beta(v^*\iota).
 \tag{D.F8}
\]
This is a proved computation for the trivial ambient bundle; it
does not assume Zhubr's unproved comparison for arbitrary
ambient tangent bundles.

Let $v_0$ be the first three vectors of $f$, and restrict to the
embedded sphere in (D.F5).  Complete a Lie tangent frame on the
sphere by its normal three-frame.  The Lie tangent frame of
$S^3$ represents an odd generator $\nu$: relative to the
bounding frame of $D^4$, its stabilized change of frame is
left quaternionic multiplication
$q\mapsto(q,qi,qj,qk)\in SO(4)$.  The associated quaternionic
line on $S^4$ has $c_2=\pm1$, hence real $p_1=\mp2$.
Its stable $SO$ class is a generator, and its $J$-image is
$\pm\nu\in\mathbb Z/24$.
If $d\in\mathbb Z/2$
is its difference from $v_0|_{S^3}$ in $\pi_3V$, then
\[
                         [S^3]_f\bmod2=1+d.
 \tag{D.F9}
\]
The quotient by normal-frame changes is exactly reduction
modulo two, by the last assertion of Zhubr's Lemma~1.
In particular $v_0$ is nontrivial relative to this sphere
precisely when the left side of (D.F9) is zero.

If $k$ is even, $v_0$ is already an admissible field in Zhubr's
Theorem~1, and its complement is trivial; hence
$\overline\gamma(X)=0$.
If $k$ is odd, $v_0$ is trivial relative to the sphere.  Choose
the admissible field $v$ supplied by that theorem.  In the
trivialization $f$, its map $v:U\to V$ has
$\langle v^*\iota,[S^3]\rangle=1$.
The exact integral Bockstein sequence and
$H^4(U;\mathbb Z)=\mathbb Z/c$ imply
\[
                         \beta(v^*\iota)=c/2.
\]
Indeed the kernel consists of reductions of integral degree-three
classes, all of which evaluate zero on the torsion sphere, while
the image is the unique order-two subgroup of $\mathbb Z/c$.
Equation (D.F8) and Zhubr's Theorem~1 now prove (D.F1).

\paragraph{Numerical consequence for the audited semistable rows.}
The independently computed integral links are $37,7,27$ for
OS55 with $(a,b)=(5,4)$ and zero indices $2,3,4$, and $4$ for
OS45 with $(a,b)=(1,8)$ and zero indices $1,2,3,4$.
Thus (D.F1) gives $\overline\gamma=2$ in the three cyclic-four
cases and $\overline\gamma=0$ in the four cyclic-eight cases.
The numerical inputs are the literal critical-cycle calculations
in \path{code/manifolds_even_critical_cycles_proof.tex}; their
independent Appell--Humbert/monopole verification belongs to the
owning D.2.44 proof.
